in the two cases under consideration. Since @@M0@@ is proper, @@M1@@ being normal (FGA 236-2.13), and since @@M2@@ is a subgroup, necessarily closed (SGA 3 VI@@M3@@ 1.9.2), for @@M4@@ to be finite it suffices that it contain no subgroup of type @@M5@@ where @@M6@@.

Consider Kummer theory, that is, the following commutative diagram, coming from the exact sequence
@@M7@@
of étale sheaves on @@M8@@ (resp. on @@M9@@):
@@M10@@
By Hilbert's ``Theorem 90'', we have @@M11@@ (resp. @@M12@@). Since @@M13@@ is proper and integral, we have @@M14@@\,; therefore @@M15@@ is injective.

It therefore suffices to see that @@M16@@ is injective, which follows from @@M17@@ being surjective, or again (SGA 1 V 6.9) from every connected étale covering @@M18@@ of @@M19@@ with group @@M20@@ having a restriction @@M21@@ which is also connected; this is a well-known consequence of Bertini's theorem and Zariski's connectedness theorem (SGA 1 X 2.10).

\sgaStableTarget{sga6.e13.rem.3-12}\textbf{Remark 3.12.} Retain the conditions of 3.11. Then:
\begin{enumerate}[label=(\roman*)]
\item The kernel @@M22@@ is a finite \uline{unipotent} group (SGA 3 XVII 1.1), since it contains no subgroup of multiplicative type @@M23@@ for any @@M24@@ (SGA 3 XVII 4.6.1 vi), which is proved in the same way, using instead
